\providecommand{\devnum}[1]{#1}

\section{Cost Prediction}
\label{sec:prediction}

The scheduler of Section~\ref{sec:scheduling} needs a number at the instant a job
arrives. This section says what that number estimates, how we keep its inputs causally
visible, which model families we compared under that one protocol, and which model we
send into the queue. All figures are development figures, and the traces they come from
are described in Section~\ref{sec:data}.

\subsection{Two Targets}

A scheduler and an evaluation want different things from the same model. The conditional
mean $\mathbb{E}[C_i \mid x_i]$ is the single-server pairwise-interchange target and is
the candidate score fitted in Section~\ref{sec:rank}; that section also explains why the
argument does not establish optimality on multiple servers, where we evaluate the order
empirically.
For comparing predictors we also need a target that stays interpretable when the mix of
jobs shifts between periods, so we score each model on a binary label. A job is
\emph{heavy} when its executed cost $C_i$ exceeds the 95th percentile of the training
rows; this is the one definition of heavy used in the paper, and the worst heavy-job
wait of Section~\ref{sec:model} is taken over the jobs it selects. On the
main judging trace that cut-off is \devnum{1.559} s and heavy jobs are \devnum{1.31\%} of
the test period; on the second judging trace it is \devnum{1{,}314} ms and
\devnum{5.07\%}. The label is deliberately fixed on training rows and never refit on the
period being scored, so its rate in the test period is a quarter of its rate in
training. Because the positive class is small we report average precision next to
AUROC, and we report RMSE on $\log(1+C)$ and Spearman correlation for the regression
output, since a ranking policy consumes order rather than level.

\subsection{The Visibility Protocol, Made Operational}

Section~\ref{sec:model} defines visibility by result availability: information about job
$j$ may enter $x_i$ exactly when $\mathrm{done}_j + \delta \le a_i$. Turning that into a
feature pipeline takes three commitments.

First, every historical feature is recomputed by streaming the log in availability order
and reading each entity's state as of the query instant. No feature is assembled over
the whole period and masked afterwards. We test this by recomputing the features from
scratch and comparing: \devnum{0} rows disagree in any of the eight configurations.
Supplementary Section~S1 gives the audit in full, including the same test applied to the
natural but wrong construction and the separate checks on the graph pipeline.

Second, thresholds are refit per target period. The heavy cut-off and the class
boundaries used inside the models are recomputed from the training rows of that target,
and each target period is predicted by a model fitted only on periods before it, frozen
before the target's first arrival. Training rows whose results would not have been
available by that instant are dropped; on the primary configuration the rule binds on
\devnum{0} rows. We report the zero because a rule that never binds is easily mistaken
for a rule that is not there.



Third, we treat the delay $\delta$ with which results reach the feature store as a
sensitivity and not as an assumption; Supplementary Section~S1 gives the two clocks in
full, on which ranking survives a ten-minute feedback delay while the probability
calibration degrades.

All families see the same rows, the same split and the same availability rule.
Supplementary Section~S1 defines them: the gradient-boosted tabular groups M1 (task
history), M2 (user history), M3 (static features of the submitted artefact), M4
(everything tabular) and M5 (M4 plus one-hop relational aggregates), with the
scrambled-neighbourhood control M6; the sequence family R1 and R1S, a GRU over the
user's own earlier jobs used alone and as a state vector inside the boosted model; and
the relational ladder N0, G0U, G0, G1 and G2, which runs a heterogeneous network from
a multilayer perceptron on M4's columns through empty neighbour slots to the real
edges and then to edges permuted inside a period and time block.

\subsection{What the Comparison Shows}

Supplementary Table~S2 gives the two judging traces. One-hop relational aggregates add
nothing on either: M5 $-$ M4 is \devnum{$+0.0019$} of AUROC on the first and
\devnum{$-0.0002$} on the second. The ladder tells the same story with a trained model in
place of hand-built aggregates. Deleting every edge (G0) and permuting the edges (G2)
leave AUROC within \devnum{0.002} of G1, and the paired intervals for G1 $-$ G0
(\devnum{$-0.0001$ $[-0.0019, 0.0018]$}) and G1 $-$ G2 (\devnum{$+0.0019$ $[-0.0000,
0.0042]$}) cover zero. The one step of the ladder that resolves is N0 $\rightarrow$ G0U,
\devnum{$+0.0139$ $[+0.0090, +0.0195]$} of AUROC at identical information, so what moves
there is the encoder, with the edge set held empty in both.

Edges do carry something, and it is confined to cold-start tasks. On the \devnum{1{,}092}
test rows whose task has fewer than five visible prior results, G1 beats M4 by
\devnum{$+0.1464$ $[+0.0997, +0.1887]$} of Spearman correlation, and beats its own
permuted-edge control by \devnum{$+0.0282$ $[+0.0130, +0.0432]$}. Over all rows the
G1 $-$ G2 Spearman difference is \devnum{$+0.0063$ $[+0.0014, +0.0111]$}, small and
resolved; the corresponding AUROC and AUPRC differences are not. We read this as a task
property rather than a verdict on relational models: when each entity's own recent record
is visible and rich, a neighbourhood adds little, and the place where it adds anything is
the entity that has no record yet.

The sequence state has the best prediction scores of anything we fitted, and no
consistent effect in the queue. R1S
reaches AUROC \devnum{0.9366} against M4's \devnum{0.9283}, a difference of
\devnum{$+0.0083$ $[+0.0038, +0.0131]$}, and it also wins on RMSE and Spearman. Inside
the queue it is not a consistent improvement: against the same ranking objective it gains
\devnum{$+0.007$ $[0.003, 0.011]$} of the deadline-window gap at busy-hour utilisation
\devnum{0.8}, loses \devnum{$-0.010$ $[-0.022, -0.002]$} at \devnum{0.5}, and is
unresolved at \devnum{1.0}. A control fitted on the same \devnum{90\%} of rows that the
GRU's early stopping left available moves the same way, so the smaller training set does
not explain it.

Which information matters is not a constant across platforms. On the first judging trace
the user's own history and the static features of the artefact beat the task's history;
on the second the task's history beats the user's by \devnum{0.11} of AUROC
(Supplementary Table~S2). A task on the first platform is attempted by one cohort within
a few days, so who is writing still carries information; a task on the second is attempted
by thousands of users over years, so the task's own record has already fixed the cost.
The non-judging traces push this further. On the
serverless trace the identity of the function explains \devnum{0.867} of the variance of
$\log(1+C)$ on its own, and static and calendar features alone are close to useless for
ranking (Spearman \devnum{0.021}). The two CI pools sit at opposite ends, entity history
giving AUROC \devnum{0.9968} on one and \devnum{0.6334} on the other, because that pool's
job labels do not separate its costs; and the compute-farm trace is predictable in rank
and not in level. Supplementary Table~S1 gives all four traces by feature group; the
AUROC column of Table~\ref{tab:cross} carries the figure the scheduling results are read
against.

\subsection{Cost of Predicting, and What Enters the Scheduler}


Prediction has to be cheap enough to sit in front of a queue whose median job is
\devnum{204} ms, and it is: Supplementary Section~S1 gives the measured per-job cost of
each family, and charging it to the predictive policies before they enqueue changes the
closed gap by less than \devnum{0.001} at every load we ran.

It is a queueing decision. M5 has the higher AUROC of the two tabular candidates and
loses in the queue at all three loads, by \devnum{0.012}, \devnum{0.008} and
\devnum{0.010} of the deadline-window gap, with intervals excluding zero in each case.
R1S has the higher AUROC of everything we fitted and does not win consistently. We
therefore take M4, the full tabular model, into Section~\ref{sec:scheduling}: it was
selected on the queueing outcome on validation periods, not on prediction quality, and
the two orderings do not agree.


% The table of predictor scores on the non-judging traces (label tab:pred_cross) moved to
% supplementary.tex as Supplementary Table S1 in the sixth revision pass, referee item G5.
% Its source comment moved with it.
%   The identity variance shares (0.867 serverless, 0.638 compute farm) quoted in the text
%     are the eta^2 column of evidence/cross_domain/out_cross_domain_table.txt; the
%     Spearman 0.021 and the two CI AUROCs 0.9968 / 0.6334 are Supplementary Table S1.
